\documentclass[11pt]{article}

\usepackage[a4paper,margin=0.78in]{geometry}
\usepackage[T1]{fontenc}
\usepackage[utf8]{inputenc}
\usepackage{lmodern}
\usepackage{amsmath,amssymb,amsthm,mathtools}
\usepackage{enumitem,booktabs,array,tabularx,microtype,xcolor}
\usepackage[numbers,sort&compress]{natbib}
\usepackage[colorlinks=true,linkcolor=blue!65!black,
  citecolor=blue!65!black,urlcolor=blue!70!black]{hyperref}
\usepackage[nameinlink,noabbrev]{cleveref}
\setlength{\emergencystretch}{3em}

\hypersetup{
  pdftitle={Source-Locator Census for the H1 Crosswalk},
  pdfauthor={Liang Wang},
  pdfsubject={Native-key collisions and production occurrence-edge census},
  pdfkeywords={twin-prime program, occurrence provenance, source locator, collision census}
}

\newtheorem{theorem}{Theorem}[section]
\newtheorem{proposition}[theorem]{Proposition}
\newtheorem{corollary}[theorem]{Corollary}
\theoremstyle{definition}
\newtheorem{definition}[theorem]{Definition}
\theoremstyle{remark}
\newtheorem{remark}[theorem]{Remark}

\newcommand{\Lzero}{\mathrm{L0}}
\newcommand{\Lone}{\mathrm{L1}}
\newcommand{\Ltwo}{\mathrm{L2}}
\newcommand{\NT}{\textnormal{\textsc{not-testable}}}
\newcommand{\FUM}{\textnormal{\textsc{frontier-unmapped}}}
\newcommand{\ETO}{\textnormal{\textsc{eligible-tail-open}}}
\newcommand{\Csrc}{\mathcal C_{\mathrm{src}}}
\newcommand{\Nkey}{\kappa_{\mathrm{nat}}}
\newcolumntype{Y}{>{\raggedright\arraybackslash}X}
\newcolumntype{P}[1]{>{\raggedright\arraybackslash}p{#1}}

\title{\textbf{Source-Locator Census for the H1 Crosswalk:}\\
\textbf{Native-Key Collisions and the Absence of Production}\\
\textbf{Occurrence Edges in the Frozen Corpus}}
\author{Liang Wang\\
School of Artificial Intelligence and Automation,\\
Huazhong University of Science and Technology, Wuhan 430074, P.R. China\\
\texttt{wangliang.f@gmail.com}}
\date{July 2026}

\begin{document}
\maketitle

\begin{abstract}
The H1 branch of the fixed-shift program currently stops at a
theorem-backed cut-to-occurrence provenance crosswalk.  We turn that
informal phrase into an auditable source-locator contract.  A positive
production edge must name a canonical source path and hash, an exact
theorem locator, an exact formula locator, and a nonempty derivation
abstract-syntax tree, while preserving the distinction between a
formal shadow and an actual occurrence.

For the frozen declared TPC-153/154/155/156/161/162 corpus, the census
contains \(2988\) distinct cut paths, two qualifying positive shadow
claims, and two formal-only scoped-obstruction claims.  It contains
no qualifying actual-occurrence edge in any of thirteen required
semantic classes.  Independently, the
archive has only \(866\) distinct native triples
\((\ell,k,d_{\rm nat})\).  Of these, \(854\) collide; \(2976\) rows lie
in collision classes, and the excess over native keys is \(2122\).
The exact multiplicity distribution is
\[
 \{1:12,\;2:220,\;4:634\}.
\]
Consequently the native triple is not a row key for this archive.

Both conclusions are deliberately scoped.  Zero qualifying edges in
the frozen corpus is not a nonexistence theorem for actual
occurrences, and failure of the native triple as an archive key is not
a failure of every enriched occurrence representation.  The next
forced structural task is to identify the smallest archived
separating key before attempting source-backed local occurrence
patches.
\end{abstract}

\noindent\textbf{Keywords:} fixed shift; occurrence provenance;
source locator; archive key; collision obstruction; proof audit.

\medskip
\noindent\fcolorbox{red!55!black}{red!3}{%
\begin{minipage}{0.96\linewidth}
\textbf{Claim boundary.}
This paper proves an exact theorem about a frozen declared source
corpus and its archive keys.  It does not prove that actual
occurrences are absent, construct a production occurrence crosswalk,
certify actual active support, select a canonical or minimal parent,
establish positive fixed-\(X\) \(\Ltwo\), give a prime-pair lower
bound, or prove the twin-prime conjecture.
\end{minipage}}

\section{Why a source census is now necessary}

TPC-153 constructs a conservative, row-separated map from every
nonsoft cut path to a \emph{partial shadow row}
\citep{WangTPC153}.  Its universal property says that any
metadata-faithful actual completion must forget back to this shadow;
it does not construct such a completion.  TPC-154 constructs two
inequivalent \emph{formal} completions and thereby stops recovery from
current artifacts alone; it explicitly leaves the augmented
theorem-backed route open \citep{WangTPC154}.  TPC-155 supplies a typed
verifier, but its production witness is absent
\citep{WangTPC155}.  TPC-156 therefore identifies
\[
 \mathsf{H1.theorem\_backed\_occurrence\_provenance\_crosswalk}
 \tag{1.1}\label{eq:first-missing}
\]
as the first missing map object \citep{WangTPC156}.  TPC-161 and
TPC-162 preserve this blocker while integrating the arithmetic return
side \citep{WangTPC161,WangTPC162}.

The phrase ``theorem-backed'' has two possible failure modes.  First,
a row may cite a file without identifying the mathematical statement
or derivation that creates the edge.  Second, a valid shadow theorem
may be silently promoted to actual-occurrence semantics.  The present
paper closes both loopholes by defining the admissible evidence unit
before counting it.

\section{Frozen corpus and evidence contract}

\begin{definition}[Frozen declared corpus]\label{def:corpus}
The corpus \(\Csrc\) consists of the canonically normalized UTF-8/LF
versions of the following eleven production artifacts:
\begin{enumerate}[leftmargin=2.1em]
 \item the TPC-153 paper, certificate, and \(2988\)-row shadow archive;
 \item the TPC-154 paper, certificate, and formal-completion archive;
 \item the TPC-155 production status and verifier audit;
 \item the TPC-156 H1 route decision;
 \item the TPC-161 occurrence--return manifest; and
 \item the TPC-162 MVP6 snapshot.
\end{enumerate}
Every source is locked by its repository-relative path and SHA-256
digest.  A digest is an integrity assertion only; it does not confer
theorem semantics.
\end{definition}

\begin{definition}[Source-locator admissibility]\label{def:admissible}
A positive production claim \(e\) is \emph{source-locator admissible}
if its record contains:
\begin{enumerate}[label=(\roman*),leftmargin=2.1em]
 \item a canonical path \(p(e)\in\Csrc\) and its recomputed canonical
       digest;
 \item an exact theorem locator that resolves in \(p(e)\);
 \item an exact formula locator that resolves in \(p(e)\);
 \item a nonempty derivation tree \(A(e)\) specifying how the cited
       formula produces the claimed typed object; and
 \item a semantic class equal to the conclusion of that derivation.
\end{enumerate}
A shadow conclusion therefore cannot inhabit an actual-occurrence
class merely by changing a status string.
\end{definition}

The derivation tree is intentionally modest.  It is not a proof
assistant kernel.  It records the operation, inputs, key fields, and
exact coefficient so that an audit can reject empty provenance and
semantic promotion.  Full mathematical truth still resides in the
cited source.

\begin{definition}[Actual H1 edge classes]\label{def:classes}
The production census checks the following thirteen classes:
\[
\begin{gathered}
\mathsf{actual\ occurrence},\quad
\mathsf{actual\ active\ support},\quad
\mathsf{canonical/minimal\ parent},\\
\mathsf{ordered\ stage},\quad
\mathsf{exact\ multiplier},\quad
Q_D,\ Q_Z,\ G,\ P_{h_0},\\
\mathsf{native\ affine\ crosswalk},\quad
\mathsf{physical\ cover},\quad
\mathsf{reconnection},\quad
\mathsf{occurrence\ registry}.
\end{gathered}
\]
These are actual-carrier classes.  A partial shadow edge is not a
member of any of them.
\end{definition}

\section{Exact source-locator census}

\begin{theorem}[Frozen-corpus edge census]\label{thm:census}
In \(\Csrc\):
\begin{enumerate}[label=(\alph*),leftmargin=2.1em]
 \item exactly two shadow claims satisfy
       \cref{def:admissible}: the TPC-153 cut-to-shadow basis injection
       and its column-conservation identity, both with semantic class
       \[
        \mathsf{PROVED\_L1\_STRUCTURAL\_SHADOW\_ONLY};
       \]
 \item exactly two TPC-154 claims satisfy the source-locator contract:
       formal completion-fibre nonuniqueness and the
       current-artifacts-only recovery obstruction, both with class
       \[
        \mathsf{FORMAL\_ONLY/SCOPED\_OBSTRUCTION};
       \]
 \item none of these four claims has actual-occurrence semantics, and
       the number of source-locator-admissible production edges in
       every class of \cref{def:classes} is zero.
\end{enumerate}
Thus the production actual-occurrence edge family is
\[
 \mathcal E_{\rm occ}(\Csrc)=\varnothing .
 \tag{3.1}\label{eq:zero-edges}
\]
\end{theorem}

\begin{proof}
The audit first recomputes all eleven path--digest locks.  It then
resolves both LaTeX labels for each positive TPC-153 claim and checks
that its derivation tree is nonempty.  The first tree is the basis
injection
\[
 e_c\longmapsto e_{\operatorname{shadow}(c)}
 \quad\text{with coefficient }1,
 \tag{3.2}\label{eq:shadow-map}
\]
and the second applies the column-sum operation to
\cref{eq:shadow-map}.  Their result types remain shadow-only.

The audit next resolves the TPC-154 completion-fibre and obstruction
theorems.  Their derivations use two conservative formal completions
with respectively one unit-weight child and two half-weight children.
The source certificate explicitly records that neither completion is
an actual-carrier lift and that only the current-artifacts-only
recovery route is stopped.  Their semantic class therefore remains
\(\mathsf{FORMAL\_ONLY/SCOPED\_OBSTRUCTION}\).

Next the audit applies an explicit per-class mapping to the frozen
production status, decision, manifest, and snapshot for the thirteen
actual classes.  TPC-155 records no
production witness, and TPC-156/161/162 all retain
\cref{eq:first-missing} as unavailable.  No record supplies the
required path, digest, two locators, derivation tree, and actual
semantic class.  Hence every class count is zero.  Mutation tests
verify that a fabricated edge, hash drift, a missing locator, an empty
tree, and a shadow-to-occurrence promotion are each rejected.
\end{proof}

\begin{remark}[Mapped census, not a future-schema scanner]
The thirteen-class table is exhaustive only for the explicitly mapped
fields of \(\Csrc\).  The program does not generically discover
arbitrary fields added by a future schema.  Any corpus or schema
extension requires a new mapping and a regenerated certificate.
\end{remark}

\begin{remark}[Closed-world scope]\label{rem:closed}
Equation \eqref{eq:zero-edges} is a closed-world statement about
\(\Csrc\).  It says
\[
 \text{``no qualifying edge is declared in this frozen corpus,''}
\]
not
\[
 \text{``no actual occurrence edge exists in mathematics.''}
\]
The latter would require a completeness theorem for the source
universe, which is neither assumed nor proved.
\end{remark}

\section{Native-key collision obstruction}

Each shadow row has a unique source-cut path identifier and carries a
native triple
\[
 \Nkey(c)=(\ell(c),k(c),d_{\rm nat}(c)).
\]
One might hope to use \(\Nkey\) as the future occurrence-row address.
The archive itself rules this out.

\begin{theorem}[Exact native collision census]\label{thm:collision}
For the \(2988\) production cut paths in \(\Csrc\), the map
\[
 \Nkey:\mathcal C\longrightarrow\mathbb Z^3
 \tag{4.1}\label{eq:native-key}
\]
has \(866\) distinct values.  Its fibre-size distribution is
\[
 \#\{v:|\Nkey^{-1}(v)|=m\}
 =
 \begin{cases}
 12,&m=1,\\
 220,&m=2,\\
 634,&m=4,\\
 0,&\text{otherwise}.
 \end{cases}
 \tag{4.2}\label{eq:distribution}
\]
Consequently:
\[
\begin{aligned}
 \#\{v:|\Nkey^{-1}(v)|>1\}&=854,\\
 \#\{c:|\Nkey^{-1}(\Nkey(c))|>1\}&=2976,\\
 \sum_v\bigl(|\Nkey^{-1}(v)|-1\bigr)&=2122.
\end{aligned}
\tag{4.3}\label{eq:collision-counts}
\]
In particular, \(\Nkey\) is not a row key for the frozen archive.
\end{theorem}

\begin{proof}
Group the \(2988\) distinct source-cut path identifiers by the exact
integer triple \eqref{eq:native-key}.  The resulting multiplicity
counter is \eqref{eq:distribution}.  The internal checks
\[
 12+220+634=866,\qquad
 12+2(220)+4(634)=2988
\]
verify the number of keys and rows.  Removing the \(12\) singleton
rows gives \(2988-12=2976\) rows in collision classes.  The excess is
\[
 220(2-1)+634(4-1)=220+1902=2122.
\]
Since a key must have singleton fibres, the final assertion follows.
\end{proof}

\begin{proposition}[Explicit fourfold collision]\label{prop:witness}
The native triple \((3,171,1)\) has four distinct archived cut paths,
distinguished by
\[
 (j_L,j_K)\in\{1,2\}\times\{7,8\}.
\]
Thus the collision in \cref{thm:collision} is not caused by duplicate
row identifiers.
\end{proposition}

\begin{proof}
The four source-cut path identifiers in the production excerpt encode
the four displayed pairs and are pairwise distinct, while their
native triples agree exactly.  The accompanying sample is copied from
the production archive and is marked as an archive excerpt, not a
synthetic or actual-occurrence witness.
\end{proof}

\section{What the two theorems jointly imply}

The two obstructions are logically different:
\[
\begin{array}{c@{\qquad}c}
 \mathcal E_{\rm occ}(\Csrc)=\varnothing
 &
 \Nkey\text{ is noninjective}
 \\[2pt]
 \text{evidence absence in a frozen corpus}
 &
 \text{exact combinatorial archive fact}.
\end{array}
\]
The first blocks an evidence-complete production lift.  The formal
TPC-154 obstruction additionally rules out canonical recovery from
the present shadow fields alone while leaving source augmentation
open.  The second
blocks a tempting shortcut for addressing its rows.  Neither prevents
an enriched construction from succeeding.

\begin{corollary}[Required enrichment]\label{cor:enrich}
Any production crosswalk that represents all \(2988\) archived cut
paths by distinct rows must use an address strictly finer than
\((\ell,k,d_{\rm nat})\), or supply an independently proved quotient
theorem identifying every colliding path.  No such quotient theorem
is present in \(\Csrc\).
\end{corollary}

\begin{proof}
This is the dichotomy between retaining distinct elements of a
noninjective fibre and identifying them by a proved equivalence
relation.  The latter is absent by the source census.
\end{proof}

\section{Reproducible certificate}

The generator reads the upstream JSON/JSONL artifacts, recomputes
canonical hashes and all counts, and serializes:
\begin{center}
\begin{tabularx}{0.96\linewidth}{@{}P{0.38\linewidth}Y@{}}
\toprule
\textbf{Artifact} & \textbf{Role}\\
\midrule
\textbf{Source census}
 & source locks, positive claims, thirteen-class census, and collision
   counts\\
\textbf{Audit}
 & recomputation checks, semantic boundaries, and mutation regressions\\
\textbf{Collision excerpt}
 & explicit fourfold production archive collision\\
\bottomrule
\end{tabularx}
\end{center}
Default mode regenerates the artifacts.  The \texttt{--check} mode
recomputes them in memory and demands byte-for-byte equality.  The
audit performs real mutations: it changes the candidate object and
then runs the same validators.  It does not merely store expected
Boolean values.

\section{Next forced object}

Before constructing local occurrence patches, the archive needs a
lossless addressing system.  TPC-163 therefore leaves one immediate
finite question:
\[
 \text{which smallest subset of existing archived fields separates
 all \(2988\) source-cut paths?}
\]
This is an exhaustive finite-key problem, not yet an actual
occurrence, active-support, or canonical-minimality theorem.  Solving
it removes ambiguity from the later local-to-global gluing interface
without overclaiming arithmetic progress.

\section{Conclusion}

TPC-163 converts the H1 crosswalk blocker into two precise facts.
The frozen declared corpus contains no theorem-backed actual
occurrence edge under an explicit source-locator contract, and the
native triple fails dramatically as an archive row key.  At the same
time, the paper preserves the correct open-world boundary: new
source-backed occurrence mathematics may still supply the missing
edges.  The next step is therefore constructive---first refine the
archive address, then formulate compatible local occurrence patches.

\bibliographystyle{plainnat}
\bibliography{references}

\end{document}
